% =====================================================================
% Topic 4 — Number Sense — part c: problems 31–40 (extending problems 1–15)
% =====================================================================
\clearpage
\dividerbar{Extension and challenge}

% ---------------------------------------------------------------- Problem 31
\begin{probbox}{31}{\relax}
Consider the numbers: $-1.5;\quad 0.25;\quad -0.75;\quad 1.2;\quad -2$.
\begin{enumerate}[label=\alph*),leftmargin=1.6em,itemsep=0pt,topsep=2pt]
  \item Arrange the numbers in increasing order.
  \item Which numbers lie between $-1$ and $1$ on the number line?
  \item Find the sum of the largest and the smallest number.
\end{enumerate}
\end{probbox}
\probsol{
\textbf{Understanding the problem.} This problem combines two skills from Problems 1 and 3: comparing decimals and comparing
negative numbers.
\begin{recallbox}Negative numbers are always less than positive numbers. Of two negative numbers, the one farther from $0$ is
smaller. For two positive numbers, compare the whole parts, then the decimal places one by one.\end{recallbox}
\step{1}{Split into two groups: negative and positive}
Negative: $-1.5;\ -0.75;\ -2$. \quad Positive: $0.25;\ 1.2$.
\step{2}{Order the negative numbers}
Distance from $0$: $-2$ is $2$ away; \ $-1.5$ is $1.5$ away; \ $-0.75$ is $0.75$ away. The one farthest from $0$ is the smallest:
\[ -2 < -1.5 < -0.75. \]
\step{3}{Order the positive numbers, then put everything together}
$0.25 < 1.2$ (whole parts $0 < 1$). Negative numbers first, positive numbers after:
\[ -2 < -1.5 < -0.75 < 0.25 < 1.2. \]
\begin{figC}\begin{scope}[x=1.6cm]
  \draw[->,thick,Navy] (-2.5,0)--(1.7,0);
  \foreach \x in {-2,-1,0,1}{\draw[Navy] (\x,0.08)--(\x,-0.08); \node[below,font=\scriptsize] at (\x,-0.1) {$\x$};}
  \foreach \x/\l in {-2/{-2},-1.5/{-1.5},-0.75/{-0.75},0.25/{0.25},1.2/{1.2}}{\fill[tmOrange] (\x,0) circle (2.6pt); \node[above,font=\scriptsize] at (\x,0.1) {$\l$};}
  \draw[tmGreen!70!black,line width=2pt,opacity=0.35] (-1,0)--(1,0);
\end{scope}\end{figC}
\step{4}{The numbers between $-1$ and $1$}
On the number line, the green segment from $-1$ to $1$ contains $-0.75$ and $0.25$. ($-1.5$ is to the left of $-1$; $1.2$ is
to the right of $1$.)
\step{5}{Sum of the largest and the smallest}
The largest is $1.2$, the smallest is $-2$:
\[ 1.2 + (-2) = -(2 - 1.2) = -0.8. \]
\begin{recallbox}Adding two numbers with opposite signs: subtract the smaller size from the larger (ignoring signs), then use
the sign of the number farther from $0$.\end{recallbox}
\begin{warnbox}Many students put $-0.75$ before $-1.5$ because ``$0.75 < 1.5$''. For negative numbers it is reversed: $-1.5 < -0.75$.\end{warnbox}
\verify{$-0.8 - 1.2 = -2$. Correct.}
\ans{a) $-2 < -1.5 < -0.75 < 0.25 < 1.2$;\quad b) $-0.75$ and $0.25$;\quad c) $-0.8$}
}

% ---------------------------------------------------------------- Problem 32
\begin{probbox}{32}{\relax}
Consider the expression $2 + 3 \times 4 - 1$.
\begin{enumerate}[label=\alph*),leftmargin=1.6em,itemsep=0pt,topsep=2pt]
  \item Evaluate the expression.
  \item Insert \emph{one} pair of brackets $(\ )$ into the expression (without changing the order of the numbers or the signs)
  to get the largest possible value. What is that largest value?
  \item Same question as b), but for the smallest value.
\end{enumerate}
\end{probbox}
\probsol{
\textbf{Understanding the problem.} This extends Problems 4 and 7: the position of the brackets changes the order of the
operations, and so changes the result.
\begin{recallbox}Brackets first; then multiplication and division; finally addition and subtraction (from left to right).\end{recallbox}
\step{1}{Value without brackets}
$2 + 3\times 4 - 1 = 2 + 12 - 1 = 13$.
\step{2}{List systematically every way to place one pair of brackets}
A pair of brackets must enclose at least two neighbouring numbers. Call the four numbers, in order, number $1, 2, 3, 4$
($2;\ 3;\ 4;\ 1$). The brackets can start at number $i$ and end at number $j$ with $i < j$: $6$ ways.
\begin{center}
\begin{tabular}{|c|l|c|}
\hline
\rowcolor{tmBlue!10} Brackets around & Expression & Value\\ \hline
numbers 1 -- 2 & $(2+3)\times 4 - 1 = 5\times 4 - 1$ & $\mathbf{19}$\\ \hline
numbers 1 -- 3 & $(2+3\times 4) - 1 = 14 - 1$ & $13$\\ \hline
numbers 1 -- 4 & $(2+3\times 4-1)$ & $13$\\ \hline
numbers 2 -- 3 & $2 + (3\times 4) - 1$ & $13$\\ \hline
numbers 2 -- 4 & $2 + (3\times 4 - 1) = 2 + 11$ & $13$\\ \hline
numbers 3 -- 4 & $2 + 3\times(4 - 1) = 2 + 9$ & $\mathbf{11}$\\ \hline
\end{tabular}
\end{center}
\step{3}{Read off the results}
The largest is $19$, with $(2+3)\times 4 - 1$: the sum $2 + 3$ is ``blown up'' $4$ times.\\
The smallest is $11$, with $2 + 3\times(4-1)$: $3$ is multiplied by a smaller number ($3$ instead of $4$).
\begin{tipbox}For a large result, let the multiplication act on a large sum; for a small result, let it act on a difference.
But to be sure, the best way is still to list every case, as in the table.\end{tipbox}
\verify{$(2+3)\times 4 - 1 = 20 - 1 = 19$ and $2 + 3\times 3 = 11$.}
\ans{a) $13$;\quad b) $(2+3)\times 4 - 1 = 19$;\quad c) $2 + 3\times(4-1) = 11$}
}

% ---------------------------------------------------------------- Problem 33
\begin{probbox}{33}{\relax}
An's storybook has $360$ pages. On the first day An reads $\dfrac{1}{4}$ of the book. On the second day An reads
$\dfrac{1}{3}$ \emph{of the remaining pages}. After two days, how many pages has An not read yet, and what fraction of the
book is that?
\end{probbox}
\probsol{
\textbf{Understanding the problem.} Careful: the fraction $\dfrac13$ on the second day is taken of the \emph{remaining
pages}, not of the whole book. This is a ``two-step'' version of Problem 9.
\begin{recallbox}$\dfrac{m}{n}$ of $a$ equals $a \div n \times m$.\end{recallbox}
\step{1}{First day}
An reads $360 \div 4 = 90$ (pages). Remaining: $360 - 90 = 270$ (pages).
\step{2}{Second day (taken of $270$ pages)}
An reads $270 \div 3 = 90$ (pages). Remaining: $270 - 90 = 180$ (pages).
\step{3}{Fraction of the book left}
$\dfrac{180}{360} = \dfrac12$.
\begin{barfigC}
  \blk{0}{2}{0}{0.6}{tmBlue}{$90$}
  \blk{2}{4}{0}{0.6}{tmGreen}{$90$}
  \blk{4}{6}{0}{0.6}{tmOrange}{$90$}
  \blk{6}{8}{0}{0.6}{tmOrange}{$90$}
  \brcup{0}{2}{0.5}{day 1}
  \brcup{2}{4}{0.5}{day 2}
  \brcup{4}{8}{0.5}{not read}
  \brcdn{2}{8}{-0.5}{left after day 1: $270$ pages, split into $3$ parts}
\end{barfigC}
\begin{warnbox}A common mistake: adding $\dfrac14 + \dfrac13 = \dfrac{7}{12}$ and saying $\dfrac{5}{12}$ of the book
($150$ pages) is left. Wrong, because $\dfrac13$ is of the \emph{remaining part} ($270$ pages), not of the whole book ($360$ pages).\end{warnbox}
\step{4}{Another way: multiply the ``remaining'' fractions}
After day 1, $\dfrac34$ of the book remains; after day 2, $\dfrac23$ of that remains, i.e.\ $\dfrac23\times\dfrac34 = \dfrac{6}{12} = \dfrac12$
of the book, which is $360 \div 2 = 180$ pages.
\verify{$90 + 90 + 180 = 360$. Correct.}
\ans{$180$ pages, which is $\dfrac12$ of the book}
}

% ---------------------------------------------------------------- Problem 34
\begin{probbox}{34}{\relax}
The first test has $40$ questions, and Bình answers $32$ of them correctly.
\begin{enumerate}[label=\alph*),leftmargin=1.6em,itemsep=0pt,topsep=2pt]
  \item What percentage of the questions in the first test did Bình answer correctly?
  \item The second test has $60$ questions. At least how many questions must Bình answer correctly in the second test so
  that, over both tests together, his correct answers make up at least $85\%$ of all the questions?
\end{enumerate}
\end{probbox}
\probsol{
\textbf{Understanding the problem.} Part a) asks for a percentage (the reverse of Problem 5). Part b) asks for a percentage of a total.
\begin{recallbox}$a$ as a percentage of $b$ is $\dfrac{a}{b}\times 100\%$. \ $p\%$ of $a$ is $a\times\dfrac{p}{100}$.\end{recallbox}
\step{1}{Percentage in the first test}
$\dfrac{32}{40}\times 100\% = 0.8\times 100\% = 80\%$.
\step{2}{Total number of questions in both tests}
$40 + 60 = 100$ (questions).
\step{3}{Correct answers needed overall}
$85\%$ of $100$ questions is $100\times\dfrac{85}{100} = 85$ (questions).
\step{4}{Correct answers needed in the second test}
Bình already has $32$ correct answers, so he needs at least $85 - 32 = 53$ (questions) in the second test.
\begin{warnbox}Do not take ``the average of the two percentages'': solving $(80\% + x) \div 2 = 85\%$ gives $x = 90\%$ of $60$
questions, i.e.\ $54$ questions -- wrong! The tests have different numbers of questions, so we must work with the
\emph{total number of questions}.\end{warnbox}
\verify{With $53$ correct: $(32 + 53) \div 100 = 85\%$ (just enough). With $52$: $84 \div 100 = 84\% < 85\%$ (not enough).}
\ans{a) $80\%$;\quad b) at least $53$ questions}
}

% ---------------------------------------------------------------- Problem 35
\begin{probbox}{35}{\relax}
A shop sells cashew nuts in packs: a $3$ kg pack costs $255$ thousand dong, a $5$ kg pack costs $410$ thousand dong (no loose sales).
\begin{enumerate}[label=\alph*),leftmargin=1.6em,itemsep=0pt,topsep=2pt]
  \item Which pack is cheaper per kilogram?
  \item Ms Hằng needs to buy \emph{exactly} $15$ kg. How should she buy them to spend as little as possible? How much does it cost?
  \item How much does it cost to buy \emph{exactly} $13$ kg?
\end{enumerate}
\end{probbox}
\probsol{
\textbf{Understanding the problem.} This extends Problem 6: compare unit prices, then find how to combine packs to get the exact mass.
\begin{recallbox}Unit price $=$ total cost $\div$ quantity.\end{recallbox}
\step{1}{Unit price of each pack}
$3$ kg pack: $255 \div 3 = 85$ (thousand dong/kg). \quad $5$ kg pack: $410 \div 5 = 82$ (thousand dong/kg).\\
So the $5$ kg pack is cheaper ($82 < 85$).
\step{2}{Ways to make exactly $15$ kg}
Let $a$ be the number of $3$ kg packs and $b$ the number of $5$ kg packs: $3a + 5b = 15$. Try $b = 0, 1, 2, 3$:
\begin{center}
\begin{tabular}{|c|c|c|c|}
\hline
\rowcolor{tmBlue!10} $b$ & $15 - 5b$ & $a$ & Cost (thousand dong)\\ \hline
$0$ & $15$ & $5$ & $5\times 255 = 1275$\\ \hline
$1$ & $10$ & not divisible by $3$ & --\\ \hline
$2$ & $5$ & not divisible by $3$ & --\\ \hline
$3$ & $0$ & $0$ & $3\times 410 = \mathbf{1230}$\\ \hline
\end{tabular}
\end{center}
Buying $3$ packs of $5$ kg is cheapest: $1\,230$ thousand dong $= 1\,230\,000$ dong.
\step{3}{Exactly $13$ kg}
$3a + 5b = 13$: \ $b = 0 \Rightarrow 13$ (not divisible by $3$); \ $b = 1 \Rightarrow 8$ (not divisible by $3$);
\ $b = 2 \Rightarrow 3$, $a = 1$. There is only one way: $1$ pack of $3$ kg and $2$ packs of $5$ kg.\\
Cost: $255 + 2\times 410 = 255 + 820 = 1075$ (thousand dong) $= 1\,075\,000$ dong.
\begin{warnbox}Part c) cannot be answered by computing $13\times 82$, because the shop does not sell loose nuts: the packs
must add up to exactly $13$ kg.\end{warnbox}
\verify{$3 + 5 + 5 = 13$ kg; $5 + 5 + 5 = 15$ kg.}
\ans{a) the $5$ kg pack ($82$ thousand dong/kg);\quad b) $3$ packs of $5$ kg, $1\,230\,000$ dong;\quad c) $1\,075\,000$ dong}
}

% ---------------------------------------------------------------- Problem 36
\begin{probbox}{36}{\relax}
The five members of a group weigh (in kg) $38;\ 41;\ 40;\ 37$ and $x$. Their mean weight is $40$ kg.
\begin{enumerate}[label=\alph*),leftmargin=1.6em,itemsep=0pt,topsep=2pt]
  \item Find $x$. Is the fifth member heavier or lighter than the mean, and by how many kilograms?
  \item Find the median weight of the group.
  \item A sixth member joins the group, and the mean weight of all six becomes $41$ kg. How much does the sixth member weigh?
\end{enumerate}
\end{probbox}
\probsol{
\textbf{Understanding the problem.} This extends Problem 14: use the sum and the deviations to find the missing number; add
the median and an ``extra member'' question.
\begin{recallbox}Sum $=$ mean $\times$ how many numbers. The deviations from the mean add up to $0$.\end{recallbox}
\step{1}{Find $x$ using deviations}
Compared with $40$: $38 \to -2$; \ $41 \to +1$; \ $40 \to 0$; \ $37 \to -3$. Sum: $-2 + 1 + 0 - 3 = -4$.\\
For the deviations to add up to $0$, the fifth member must be $+4$: $x = 40 + 4 = 44$ (kg), $4$ kg heavier than the mean.\\
Check with the sum: $5\times 40 = 200$; \ $200 - (38 + 41 + 40 + 37) = 200 - 156 = 44$.
\step{2}{Median}
In increasing order: $37;\ 38;\ \mathbf{40};\ 41;\ 44$. The middle (third) number is $40$. The median is $40$ kg.
\step{3}{The sixth member}
Total weight of six members: $6\times 41 = 246$ (kg). The first five weigh $200$ kg in total. The sixth weighs $246 - 200 = 46$ (kg).
\begin{tipbox}Think in deviations: the mean rises by $1$ kg for all $6$ people, so $6$ kg extra above the level of $40$ are needed;
the new member must weigh $40 + 6 = 46$ kg.\end{tipbox}
\begin{warnbox}It is not ``the sixth member weighs $41$ kg'' -- then the mean would rise only a little: $241 \div 6 \approx 40.17$ kg.\end{warnbox}
\verify{$(37 + 38 + 40 + 41 + 44 + 46) \div 6 = 246 \div 6 = 41$. Correct.}
\ans{a) $x = 44$ kg, $4$ kg heavier than the mean;\quad b) $40$ kg;\quad c) $46$ kg}
}

% ---------------------------------------------------------------- Problem 37
\begin{probbox}{37}{\stars{1}}
A running track around the schoolyard is $200$ m long. Mai runs $6$ laps in $4$ minutes. Tú runs at a constant $18$ km/h.
\begin{enumerate}[label=\alph*),leftmargin=1.6em,itemsep=0pt,topsep=2pt]
  \item Who is faster?
  \item Keeping the same speeds, how many minutes does each of them need to run $3$ km?
  \item They start at the same time from the same point of the track, but run in \emph{opposite directions}. After how many
  seconds do they meet for the first time?
  \item One day Tú increases his speed to $21.6$ km/h while Mai runs as before. They start at the same time from the same
  point, running in the \emph{same direction}. How long is it until Tú is exactly one lap ahead of Mai?
\end{enumerate}
\end{probbox}
\probsol{
\textbf{Understanding the problem.} This extends Problem 12: find a speed from laps run, convert km/h to m/s, then study two
runners on the same circular track.
\begin{recallbox}$v = s \div t$, \ $t = s \div v$. \ $1$ km/h $= \dfrac{5}{18}$ m/s, so to convert km/h to m/s multiply by $5$ and divide by $18$.\end{recallbox}
\step{1}{Mai's speed}
Distance: $6\times 200 = 1200$ (m). Time: $4$ minutes $= 240$ seconds. \ $v_{\text{Mai}} = 1200 \div 240 = 5$ (m/s).
\step{2}{Tú's speed (in m/s)}
$v_{\text{Tú}} = 18\times\dfrac{5}{18} = 5$ (m/s). They run \textbf{equally fast}!
\begin{warnbox}Comparing ``$5$'' with ``$18$'' directly is meaningless, because the units differ (m/s and km/h). Convert to the
same unit first.\end{warnbox}
\step{3}{Time to run $3$ km}
$3$ km $= 3000$ m. \ $t = 3000 \div 5 = 600$ (seconds) $= 600 \div 60 = 10$ (minutes). Both need $10$ minutes.
\step{4}{Opposite directions}
Every second Mai runs $5$ m and Tú runs $5$ m the other way, so the ``combined'' distance they cover grows by
$5 + 5 = 10$ m per second. They first meet when the \emph{total} of their distances is exactly one lap, $200$ m:
\[ t = 200 \div 10 = 20 \text{ (seconds)}. \]
\begin{figC}
  \draw[Slate!60,line width=4pt] (0,0) circle (1.1);
  \fill[Navy] (0,1.1) circle (2.5pt); \node[above=4pt,font=\scriptsize] at (0,1.35) {start};
  \draw[->,tmBlue,thick] (95:1.4) arc (95:265:1.4);
  \draw[->,tmOrange,thick] (85:1.4) arc (85:-85:1.4);
  \node[left,font=\scriptsize,text=tmBlue!80!black] at (-1.5,0) {Mai: $100$ m};
  \node[right,font=\scriptsize,text=tmOrange!80!black] at (1.5,0) {Tú: $100$ m};
  \fill[AnsRed] (0,-1.1) circle (2.5pt); \node[below,font=\scriptsize,text=AnsRed] at (0,-1.5) {they meet after $20$ s};
\end{figC}
After $20$ seconds each has run $5\times 20 = 100$ m, exactly half a lap: they meet at the point opposite the start.
\step{5}{Same direction, Tú faster}
$v_{\text{Tú}} = 21.6\times\dfrac{5}{18} = \dfrac{108}{18} = 6$ (m/s). Every second Tú runs $6 - 5 = 1$ m more than Mai. Tú is
exactly one lap ahead when he has run $200$ m more than Mai:
\[ t = 200 \div 1 = 200 \text{ (seconds)} = 3 \text{ minutes } 20 \text{ seconds}. \]
\verify{After $200$ seconds: Tú has run $6\times 200 = 1200$ m ($6$ laps), Mai $5\times 200 = 1000$ m ($5$ laps). Tú is exactly $1$ lap ahead.}
\begin{parentbox}You can act out parts c) and d) with two coins on a round plate: in opposite directions the ``gap to close''
shrinks by the \emph{sum} of the speeds; in the same direction it shrinks by the \emph{difference} of the speeds.\end{parentbox}
\ans{a) equally fast ($5$ m/s);\quad b) $10$ minutes;\quad c) $20$ seconds;\quad d) $200$ seconds ($3$ minutes $20$ seconds)}
}

% ---------------------------------------------------------------- Problem 38
\begin{probbox}{38}{\stars{1}}
\begin{enumerate}[label=\alph*),leftmargin=1.6em,itemsep=0pt,topsep=0pt]
  \item Find the digits $a$, $b$ such that the number $\nd{4a7b}$ is divisible by $2$, $5$ and $9$.
  \item How many four-digit numbers of the form $\nd{5a3b}$ are divisible by $9$? List them.
\end{enumerate}
\end{probbox}
\probsol{
\textbf{Understanding the problem.} This extends Problem 10: use several divisibility rules at once and count systematically.
\begin{recallbox}Divisible by $2$ and $5$ $\iff$ the last digit is $0$. \ Divisible by $9$ $\iff$ the digit sum is divisible by $9$.\end{recallbox}
\begin{proofbox}Why is ``divisible by both $2$ and $5$'' the same as ``ends in $0$''? Divisibility by $2$ needs an even last digit;
divisibility by $5$ needs a last digit of $0$ or $5$. The only digit that is even and in $\{0;5\}$ is $0$.\end{proofbox}
\step{1}{Part a) -- find $b$}
$\nd{4a7b}$ is divisible by $2$ and $5$, so $b = 0$.
\step{2}{Part a) -- find $a$}
Digit sum: $4 + a + 7 + 0 = 11 + a$. Since $0 \le a \le 9$, we have $11 \le 11 + a \le 20$. The only multiple of $9$ in this
range is $18$. So $11 + a = 18$, $a = 7$. The number is $4770$.
\verify{$4770 = 9\times 530$, and it ends in $0$, so it is divisible by $2$ and $5$.}
\step{3}{Part b) -- the condition}
The digit sum of $\nd{5a3b}$ is $5 + a + 3 + b = 8 + a + b$. Since $0 \le a + b \le 18$, we have $8 \le 8 + a + b \le 26$. The
multiples of $9$ in this range are $9$ and $18$. Hence $a + b = 1$ or $a + b = 10$.
\step{4}{Part b) -- count systematically}
$a + b = 1$: $(a;b) = (0;1)$ or $(1;0)$ \ -- \ $2$ numbers.\\
$a + b = 10$: $a$ runs from $1$ to $9$, $b = 10 - a$ (from $9$ down to $1$) \ -- \ $9$ numbers. ($a = 0$ would give $b = 10$, not a digit.)\\
In total: $2 + 9 = 11$ numbers:
\[ 5031,\ 5130,\ 5139,\ 5238,\ 5337,\ 5436,\ 5535,\ 5634,\ 5733,\ 5832,\ 5931. \]
\begin{warnbox}Do not forget that $a$ may be $0$ (the first digit is $5$), and do not miss the case where the digit sum is $9$
(that is, $a + b = 1$).\end{warnbox}
\verify{Each number above has digit sum $9$ or $18$, e.g.\ $5\,436$: $5+4+3+6 = 18$.}
\ans{a) $a = 7$, $b = 0$ (the number $4770$);\quad b) $11$ numbers}
}

% ---------------------------------------------------------------- Problem 39
\begin{probbox}{39}{\stars{1}}
In the addition below, $M$ and $N$ are two different non-zero digits. Find all pairs $(M;N)$ that work.
\[ \begin{array}{r@{\;}c@{\;}c} & M & N\\ & N & M\\ + & N & N\\ \hline 1 & 8 & 7 \end{array} \]
\end{probbox}
\probsol{
\textbf{Understanding the problem.} This extends Problem 13. This time, instead of trying column by column, we use place value
to get one simple equation, and then find \emph{all} solutions.
\begin{recallbox}$\nd{MN} = 10M + N$, \ $\nd{NM} = 10N + M$, \ $\nd{NN} = 10N + N = 11N$.\end{recallbox}
\step{1}{Write the sum using place value}
\[ \nd{MN} + \nd{NM} + \nd{NN} = (10M + N) + (10N + M) + 11N = 11M + 22N = 11\times(M + 2N). \]
\step{2}{Set up the equation}
$11\times(M + 2N) = 187$, so $M + 2N = 187 \div 11 = 17$.
\step{3}{Find every pair of digits}
$M = 17 - 2N$ must be a digit from $1$ to $9$, so $17 - 2N \le 9$, i.e.\ $N \ge 4$; and $17 - 2N \ge 1$, i.e.\ $N \le 8$.
\begin{center}
\begin{tabular}{|c|c|c|c|c|c|}
\hline
\rowcolor{tmBlue!10} $N$ & $4$ & $5$ & $6$ & $7$ & $8$\\ \hline
$M = 17 - 2N$ & $9$ & $7$ & $5$ & $3$ & $1$\\ \hline
Addition & $94+49+44$ & $75+57+55$ & $56+65+66$ & $37+73+77$ & $18+81+88$\\ \hline
\end{tabular}
\end{center}
All five pairs have $M \ne N$, and every sum equals $187$.
\begin{tipbox}When letters appear ``symmetrically'', as in $\nd{MN}$ and $\nd{NM}$, try place value: there is often a common
factor ($11$ here) that makes the problem much simpler.\end{tipbox}
\begin{warnbox}The problem asks for \emph{all} pairs. Finding one pair (say $M = 9$, $N = 4$) and stopping is not enough.\end{warnbox}
\verify{$94 + 49 + 44 = 187$; \ $75 + 57 + 55 = 187$; \ $56 + 65 + 66 = 187$; \ $37 + 73 + 77 = 187$; \ $18 + 81 + 88 = 187$.}
\ans{$(M;N) \in \{(9;4),\ (7;5),\ (5;6),\ (3;7),\ (1;8)\}$}
}

% ---------------------------------------------------------------- Problem 40
\begin{probbox}{40}{\stars{2}}
A sequence is built by the Fibonacci rule: from the third number on, each number is the sum of the two numbers just before it.
\begin{enumerate}[label=\alph*),leftmargin=1.6em,itemsep=0pt,topsep=2pt]
  \item The first number is $2$ and the fifth number is $19$. Find the second and the eighth numbers of the sequence.
  \item Find the mean of the first five numbers of the sequence in part a).
  \item In the familiar Fibonacci sequence $1, 1, 2, 3, 5, 8, \ldots$, is the $2026$th number even or odd? Explain.
\end{enumerate}
\end{probbox}
\probsol{
\textbf{Understanding the problem.} This extends Problem 15. In part a) the second number is unknown, so we call it $b$ and
express the following numbers in terms of $b$. In part c) we cannot write out $2026$ numbers; we need a pattern of even and odd.
\step{1}{Express the numbers in terms of $b$}
\begin{center}
\begin{tabular}{|c|c|c|c|c|c|}
\hline
\rowcolor{tmBlue!10} Position & 1 & 2 & 3 & 4 & 5\\ \hline
Number & $2$ & $b$ & $2 + b$ & $b + (2 + b) = 2 + 2b$ & $(2 + b) + (2 + 2b) = 4 + 3b$\\ \hline
\end{tabular}
\end{center}
\step{2}{Find $b$}
$4 + 3b = 19 \Rightarrow 3b = 15 \Rightarrow b = 5$. The second number is $5$.
\step{3}{Continue up to the eighth number}
$2,\ 5,\ 7,\ 12,\ 19,\ 31,\ 50,\ 81$ \quad ($12 + 19 = 31$; \ $19 + 31 = 50$; \ $31 + 50 = 81$). The eighth number is $81$.
\step{4}{Mean of the first five numbers}
$(2 + 5 + 7 + 12 + 19) \div 5 = 45 \div 5 = 9$.
\step{5}{The even--odd pattern of the Fibonacci sequence}
\begin{recallbox}odd $+$ odd $=$ even; \ odd $+$ even $=$ odd; \ even $+$ odd $=$ odd.\end{recallbox}
Write O (odd) and E (even) for the numbers $1, 1, 2, 3, 5, 8, 13, 21, 34, \ldots$:
\begin{center}
\begin{tabular}{|c|c|c|c|c|c|c|c|c|c|}
\hline
\rowcolor{tmBlue!10} Position & 1 & 2 & 3 & 4 & 5 & 6 & 7 & 8 & 9\\ \hline
Number & 1 & 1 & 2 & 3 & 5 & 8 & 13 & 21 & 34\\ \hline
Even/odd & O & O & \textbf{E} & O & O & \textbf{E} & O & O & \textbf{E}\\ \hline
\end{tabular}
\end{center}
The pattern ``O, O, E'' repeats forever.
\begin{proofbox}Suppose at some point we have two consecutive numbers O, O. Then the next number is O $+$ O $=$ E; the one after is
O $+$ E $=$ O; the one after that is E $+$ O $=$ O. We again have two consecutive numbers O, O, so the cycle ``O, O, E'' starts
again. The sequence begins with $1, 1$ (O, O), so this pattern repeats all the way: the numbers in positions divisible by $3$
are even, all the others are odd.\end{proofbox}
\step{6}{Look at position $2026$}
$2026 = 3\times 675 + 1$, i.e.\ $2026$ leaves remainder $1$ when divided by $3$. This position corresponds to the first ``O'' in
a cycle. So the $2026$th number is \textbf{odd}.
\begin{warnbox}You cannot conclude ``$2026$ is even, so the $2026$th number is even too''. The parity of the \emph{position} is
different from the parity of the \emph{number} in that position; for example, the $2$nd number is $1$ (odd).\end{warnbox}
\verify{Part a): $2 + 5 = 7$, $5 + 7 = 12$, $7 + 12 = 19$ -- indeed the fifth number is $19$. Part c): test the pattern on small
positions: position $4 = 3\times 1 + 1$ holds $3$ (odd); position $7 = 3\times 2 + 1$ holds $13$ (odd).}
\ans{a) the second number is $5$, the eighth is $81$;\quad b) $9$;\quad c) odd}
}
